\chapter{The prototype so far}\label{ch:prototype}
% Draft by Claude from our old iteration notes, edit freely.
% This chapter is a snapshot up to a milestone, not a list of iterations. Keep adding to it
% until the semester ends. For each part: what we did, and WHY.

This chapter describes the prototype as it is on \fillin{milestone date}. Instead of listing
every version, each section describes one part of the device: what it does now, and why we
made it that way. Details of each printed version (photos, notes) are kept in the
\texttt{iterations/} folder of the project repository.

\section{How the device works}\label{sec:how-it-works}

The device is a small, automated version of the Hirst trap described in
Chapter~\ref{ch:introduction}:

\begin{enumerate}
  \item A fan moves outdoor air onto a strip of clear sticky tape, so particles stick to it.
  \item A stepper motor moves the tape on, so the sampled piece lies over the camera.
  \item An LED shines through the tape, and the camera underneath takes a picture.
  \item The picture is saved with the time and all settings. Later it is sent to a server.
  \item This repeats. The used tape is kept on a roll, so it can be checked under a microscope later.
\end{enumerate}

A Raspberry Pi~3~B+ controls everything. \fillin{Photo of the current prototype.}

\section{Enclosure and tape transport}\label{sec:enclosure}

We started from the earlier design (Section~\ref{sec:research-question}), but drew our own 3D
model from scratch. The reasons were:

\begin{itemize}
  \item \textbf{Easy to build.} The goal is that anyone can build one, so we avoided screws
  and other extra parts where we could. Parts hold by friction or slot into each other.
  \item \textbf{One print.} We have access to a larger 3D printer, so parts can be bigger and
  fewer.
\end{itemize}

\paragraph{LED tower.}
% Draft by Claude 2026-11-05, edit freely.
The LED sits in a black printed tower that stands on the camera box, so it is always at the same
height above the tape and no room light reaches the sensor. The light shines through a small
hole at the bottom of the tower, because a small light source gives sharper shadows
(Section~\ref{sec:imaging-test-3}). The tape runs under the tower, between its legs, so nothing
touches the sticky side. The tower snaps onto the camera box and hooks over the top of the frame.
It can be pulled off and put back without taking the tape out, which we need often while testing.
Printing it in a way so that there is only a pinhole gap 0.6mm between the light source and the
sensor, which is the smallest our printer is able do to, does seem to give better results.

\paragraph{Blowing or sucking.}
The reference trap sucks air through a narrow slit onto the tape. That gives a known airflow
and a sharp, narrow line of particles, which is needed to report pollen per cubic metre of air.
For the first prototype we chose to blow air onto the tape instead, because it is much simpler
to build. 

\section{Electronics}\label{sec:electronics}

\begin{itemize}
  \item \textbf{Stepper motor.} A 28BYJ-48 with a ULN2003 driver board moves the tape. It is
  slow and weak, but cheap and strong enough for tape.
  \item \textbf{Fan.} A 5~V fan, switched by a MOSFET, because a Raspberry Pi pin cannot power a
  fan directly. A diode protects the MOSFET.
  \item \textbf{Camera.} The Raspberry Pi HQ camera with the lens removed
  (Section~\ref{sec:imaging-tests}).
  \item \textbf{LED.} One LED cut from a string of colour Christmas lights we had at home.
  See the light source below.
\end{itemize}

\paragraph{Light source.}
% Draft by Claude 2026-11-07, edit freely.
We changed the light twice. Each time the pictures told us why.
\begin{enumerate}
  \item \textbf{Blue 5~mm LED} (Test~1). Pure blue light only reaches one in four pixels,
  because of the colour filter on the sensor, so the pictures were dark and noisy.
  \item \textbf{RGB 5~mm LED} (Test~2). Each colour has its own Pi pin, so the software can
  choose the colour. But it was dim. A Pi pin gives only 3.3~V, and the green and blue parts
  need most of that just to light up, so very little current is left.
  \item \textbf{Christmas light} (Test~3). One LED cut from a chain of Christmas
  lights. The nice part is that the LED only needs power
  from the 5~V supply, ground, and one data wire from the Pi. It is much brighter than the
  RGB LED. The LED is also surrounded in a slightly cloudy, rubbery material that difuses the
  light. This mixes the three colours a little bit before they reach the sensor, so a mixed colour such as
  white seems to cast one shadow instead of two or three shifted shadows.
  \future{Check this with a white-light picture.}
  We used it because we had it. A future kit should use a proper part of the same kind (an
  addressable RGB LED, such as the WS2812, with a diffuser), which is cheap and easy to buy.
\end{enumerate}
